\chapter{Lesson Learned}

\section{Problems Faced During this Period}
Designing, developing, and deploying an integrated automotive enterprise management platform during my internship at Radiant Pharmaceuticals Limited presented several significant software engineering challenges. The primary technical hurdle involved database aggregation bottlenecks within Django's default Object-Relational Mapping (ORM) layer. Because the star-schema analytical model required multi-table joins across sales facts, physical stores, vehicle specifications, customers, and employees, standard ORM querysets generated substantial Python object hydration overhead. This resulted in noticeable processing delays and heavy memory consumption during complex reporting routines.

Furthermore, bridging the public e-commerce storefront with the internal dealership ERP engine introduced architectural coupling and transactional concurrency risks. Because physical vehicles represent unique single-unit inventory items, simultaneous reservation attempts by online buyers created potential race conditions where multiple customers could attempt to check out the same stock unit. Managing the corporate supervisory hierarchy also proved challenging, as navigating deep multi-tier reporting lines through recursive relational queries degraded permission verification speeds across dealership branches. Finally, handling high-resolution vehicle media and icon web fonts across differing deployment environments initially triggered browser security restrictions and layout shifts.

\section{Solution of those Problems}
To resolve the database performance bottlenecks, I bypassed ORM model hydration by implementing optimized raw SQL aggregation queries executed directly through Django's database cursor interface. By utilizing hand-crafted multi-table joins and composite grouping clauses on indexed foreign keys, analytical query execution became near-instantaneous, allowing executive dashboards to render seamlessly without placing excessive load on the server.

To eliminate race conditions and ensure data consistency, all critical checkout transitions—including vehicle status updates, holding order creations, and invoice generations—were encapsulated within atomic database transaction blocks using Django's \texttt{@transaction.atomic} decorator alongside row-level locking via \texttt{select\_for\_update()}. Architectural decoupling was maintained by routing inter-application business logic through dedicated service layers rather than direct model dependencies. In addition, supervisory hierarchy traversal was streamlined by introducing a flat hierarchy mapping model that stores direct references across management tiers, enabling constant-time permission checks. Asset delivery was stabilized by serving fonts and scripts through a same-origin static pipeline while applying responsive CSS wrappers and lazy-loading attributes to catalog media.

\section{Professional and Corporate Lessons Learned}
Beyond technical problem solving, working within the Management Information Systems (MIS) department at Radiant Pharmaceuticals Limited provided vital insights into enterprise software practices. I learned how corporate IT divisions prioritize data consistency, auditability, and clear reporting interfaces to support strategic executive decisions. The experience reinforced the necessity of profiling and empirical verification rather than relying on theoretical design assumptions. Collaborating closely with industry supervisor Ayan Chowdhury and academic supervisor Azwad Abid taught me how to align iterative development milestones with both business requirements and academic engineering standards.